\documentclass[11pt]{article}

\usepackage[T1]{fontenc}
\usepackage{lmodern}
\usepackage{microtype}
\usepackage[margin=0.9in]{geometry}
\usepackage{amsmath,amssymb,amsthm,mathtools}
\usepackage{booktabs,longtable,array}
\usepackage{enumitem}
\usepackage{xcolor}
\usepackage{hyperref}
\usepackage[nameinlink,noabbrev]{cleveref}
\usepackage{fancyhdr}
\usepackage{lastpage}

\hypersetup{
  colorlinks=true,
  linkcolor=blue!55!black,
  citecolor=blue!55!black,
  urlcolor=blue!55!black,
  pdftitle={A complete candidate proof of Erdos Problem 971},
  pdfauthor={KyungMin Han}
}

\pagestyle{fancy}
\fancyhf{}
\lhead{Candidate proof of Erd\H{o}s Problem 971}
\rhead{Page \thepage\ of \pageref{LastPage}}
\setlength{\headheight}{14pt}

\newtheorem{theorem}{Theorem}[section]
\newtheorem{proposition}[theorem]{Proposition}
\newtheorem{lemma}[theorem]{Lemma}
\newtheorem{corollary}[theorem]{Corollary}
\theoremstyle{definition}
\newtheorem{definition}[theorem]{Definition}
\theoremstyle{remark}
\newtheorem{remark}[theorem]{Remark}
\newtheorem{checkpoint}[theorem]{Integrity checkpoint}

\newcommand{\ph}{\varphi}
\newcommand{\Ss}{\mathfrak S}
\newcommand{\Pp}{\mathbb P}
\newcommand{\N}{\mathbb N}
\newcommand{\R}{\mathbb R}
\newcommand{\eps}{\varepsilon}
\newcommand{\1}{\mathbf 1}
\newcommand{\OO}{\mathrm O}
\newcommand{\om}{\omega}

\definecolor{statusblue}{RGB}{235,244,252}
\definecolor{statusborder}{RGB}{48,97,145}
\definecolor{warningbg}{RGB}{255,248,226}
\definecolor{warningborder}{RGB}{164,116,20}
\definecolor{passbg}{RGB}{237,249,239}
\definecolor{passborder}{RGB}{45,120,59}

\title{\bfseries A Complete Candidate Proof of Erd\H{o}s Problem 971\\[0.4em]
\large Pointwise variance and a uniform three-tuple upper sieve}
\author{KyungMin Han}
\date{26 July 2026}

\begin{document}
\maketitle



\begin{center}
\fcolorbox{warningborder}{warningbg}{%
\begin{minipage}{0.92\textwidth}
\textbf{AI-assistance disclosure.}
The research exploration, proof organization, drafting, and formalization workflow for this manuscript were substantially assisted using OpenAI's GPT-5.6 Pro.
The author has reviewed the argument and assumes responsibility for every mathematical claim, citation, and error in the manuscript.
\end{minipage}}
\end{center}

\begin{abstract}
For a reduced residue class $a\pmod q$, let $p(a,q)$ be the least prime congruent to $a\pmod q$.
We prove, subject only to the correctness of cited published theorems, that there are absolute constants $c,C>0$ such that, for every sufficiently large $q$,
\[
 \#\{a\pmod q:(a,q)=1,\ p(a,q)>(1+c)\ph(q)\log q\}\ge C\ph(q).
\]
At the critical cutoff $X\sim\ph(q)\log q$, the pointwise Friedlander--Goldston variance lower bound forces the second factorial moment of the prime occupancies to satisfy $C_2\gg\ph(q)$.
A uniform three-linear-form upper sieve, followed by an explicit average of the associated singular series, gives $C_3\ll\ph(q)$ even for very smooth moduli.
These two bounds force a positive proportion of residue classes to contain at least two primes below $X$.
Since the total number of prime incidences is only $\ph(q)+o(\ph(q))$, a positive proportion of classes are empty.
The prime number theorem then preserves a positive proportion of empty classes after the cutoff is enlarged to $(1+c)\ph(q)\log q$.
Every previously identified issue---centering, floors, prime powers, nonadmissible triples, coefficient uniformity, and smooth-modulus local factors---is treated explicitly.
\end{abstract}

\tableofcontents

\section{Statement and external inputs}

For $(a,q)=1$, define
\[
 p(a,q):=\min\{p\in\Pp:p\equiv a\pmod q\}.
\]
Dirichlet's theorem guarantees that this minimum exists.
The problem asks whether there is an absolute $c>0$ such that, for every sufficiently large $q$,
\begin{equation}\label{eq:target}
 \#\{a\pmod q:(a,q)=1,\ p(a,q)>(1+c)\ph(q)\log q\}\gg\ph(q).
\end{equation}

The proof uses the following two non-elementary theorems.

\begin{theorem}[Friedlander--Goldston pointwise variance lower bound]\label{thm:FG}
Let
\[
 V(x;q):=\sum_{a\bmod q}^{*}
 \left(\psi(x;q,a)-\frac{\psi(x,\chi_0)}{\ph(q)}\right)^2.
\]
For every fixed $A>0$, uniformly for
\[
 \frac{x}{(\log x)^A}\le q\le x,
\]
one has
\[
 V(x;q)\ge\left(\frac12-o(1)\right)x\log q
\]
as $x\to\infty$.
\end{theorem}

This formulation is recorded by Fiorilli from Friedlander and Goldston's theorem.  The original paper uses an equivalent centering by $x/\ph(q)$; the transfer between the two centerings is checked in \cref{sec:pair}.

\begin{theorem}[Uniform three-tuple upper bound]\label{thm:Dubbe}
Let $a_1,a_2,b_1,b_2\in\mathbb Z$, and suppose the three forms
\[
 m,\qquad a_1m+b_1,\qquad a_2m+b_2
\]
are nondegenerate and admissible.  Then, uniformly in the coefficients, for all sufficiently large $x$,
\[
 \#\{p\le x:p,\ a_1p+b_1,\ a_2p+b_2\text{ are prime}\}
 \ll \Ss\,\frac{x}{(\log x)^3},
\]
where
\[
 \Ss=\prod_{\ell}
 \left(1-\frac{\rho(\ell)}{\ell}\right)
 \left(1-\frac1\ell\right)^{-3}
\]
and $\rho(\ell)$ is the number of roots modulo $\ell$ of the product of the three forms.
The implied constant is absolute.
\end{theorem}

For direct reproducibility one may use Dubbe's explicit Theorem 1.  In the case of three forms it gives the factor $48$ times the singular series and an explicit correction tending to $1$.  For example, once $\log x\ge100$, its polynomial correction is less than $2$, so an absolute constant $96$ is sufficient.

\begin{remark}
Both inputs are unconditional.  The proof below does not use GRH, Elliott--Halberstam, Hardy--Littlewood, or a lower-bound sieve for a fixed prime pair.
\end{remark}

\section{The critical cutoff}

Write
\[
 \Phi:=\ph(q),\qquad
 R:=\left\lfloor\frac{\Phi\log q}{q}\right\rfloor,
 \qquad X:=qR.
\]
The standard uniform estimate
\[
 \frac{q}{\ph(q)}\ll\log\log q
\]
implies
\[
 R\gg\frac{\log q}{\log\log q}\longrightarrow\infty.
\]
Hence $R\ge1$ and $X\ge q$ for all sufficiently large $q$.
Moreover,
\begin{equation}\label{eq:floor}
 0\le\Phi\log q-X<q,
\end{equation}
so
\begin{equation}\label{eq:Xasymp}
 X=(1+o(1))\Phi\log q.
\end{equation}
Since $R\le\log q$,
\begin{equation}\label{eq:FG-range}
 q=\frac XR\ge\frac{X}{\log q}>\frac{X}{\log X}.
\end{equation}
Thus $(X,q)$ lies in the range of \cref{thm:FG} with $A=1$.
We also have
\begin{align}
 \frac X\Phi
 &=\log q+\OO\!\left(\frac q\Phi\right)
  =\log q+\OO(\log\log q),\label{eq:Xoverphi}\\
 \log X&=\log q+\log R=\log q+\OO(\log\log q),\label{eq:logX}\\
 \frac{X}{\log X}&=(1+o(1))\Phi.\label{eq:XlogX}
\end{align}

For each reduced residue class $a\pmod q$, put
\[
 N_a=N_a(X):=\#\{p\le X:p\equiv a\pmod q\}
\]
and define
\[
 C_2:=\sum_{a\bmod q}^{*}\binom{N_a}{2},
 \qquad
 C_3:=\sum_{a\bmod q}^{*}\binom{N_a}{3}.
\]
We prove
\begin{equation}\label{eq:moments-goal}
 C_2\gg\Phi,\qquad C_3\ll\Phi.
\end{equation}

\section{The second factorial moment from pointwise variance}\label{sec:pair}

For a reduced residue class $a\pmod q$, let
\[
 \psi_a:=\sum_{\substack{n\le X\\n\equiv a\ (q)}}\Lambda(n),
 \qquad
 S:=\sum_{\substack{n\le X\\(n,q)=1}}\Lambda(n)=\psi(X,\chi_0).
\]
Set
\[
 V_S(X;q):=\sum_{a\bmod q}^{*}
 \left(\psi_a-\frac S\Phi\right)^2.
\]

\subsection{Centering and exact expansion}

If the cited variance theorem is stated with center $X/\Phi$, then
\begin{equation}\label{eq:center-transfer}
 \sum_a^*\left(\psi_a-\frac X\Phi\right)^2
 =V_S(X;q)+\frac{(S-X)^2}{\Phi}.
\end{equation}
The prime number theorem gives $S=X+o(X)$, uniformly along the present sequence of $q$, because the omitted prime powers based on primes dividing $q$ contribute at most
\[
 \sum_{p\mid q}\left\lfloor\frac{\log X}{\log p}\right\rfloor\log p
 \le \om(q)\log X\ll(\log q)(\log X)=o(X).
\]
Consequently
\[
 \frac{(S-X)^2}{\Phi}=o\!\left(\frac{X^2}{\Phi}\right)
 =o(X\log q),
\]
by \cref{eq:Xoverphi}.  Thus the two centerings have the same lower bound for our purpose.

Define
\[
 D_\Lambda:=\sum_{\substack{n\le X\\(n,q)=1}}\Lambda(n)^2
\]
and
\[
 P_\Lambda:=
 \sum_{\substack{m<n\le X\\m\equiv n\ (q)\\(mn,q)=1}}
 \Lambda(m)\Lambda(n).
\]
Opening the square gives the exact identity
\begin{equation}\label{eq:variance-expand}
 V_S(X;q)=D_\Lambda+2P_\Lambda-\frac{S^2}{\Phi}.
\end{equation}

\subsection{Diagonal cancellation}

The prime number theorem and partial summation give
\begin{equation}\label{eq:D-asymp}
 \sum_{n\le X}\Lambda(n)^2=X\log X+\OO(X).
\end{equation}
Removing the powers of primes dividing $q$ costs at most
\[
 \sum_{p\mid q}
 \left\lfloor\frac{\log X}{\log p}\right\rfloor(\log p)^2
 \le (\log X)\sum_{p\mid q}\log p
 \le(\log X)\log q=o(X\log q).
\]
Therefore
\begin{equation}\label{eq:D-coprime}
 D_\Lambda=X\log X+\OO(X)+o(X\log q).
\end{equation}
Also $S=X+o(X)$, so
\[
 \frac{S^2}{\Phi}=\frac{X^2}{\Phi}+o(X\log q).
\]
Using \cref{eq:Xoverphi,eq:logX},
\begin{align}
 D_\Lambda-\frac{S^2}{\Phi}
 &=X\left(\log X-\frac X\Phi\right)+o(X\log q)\\
 &=\OO(X\log\log q)+o(X\log q)\\
 &=o(X\log q).\label{eq:diag-cancel}
\end{align}

Combining \cref{thm:FG,eq:variance-expand,eq:diag-cancel},
\begin{equation}\label{eq:weighted-pairs}
 P_\Lambda\ge\left(\frac14-o(1)\right)X\log q.
\end{equation}

\subsection{Removal of proper prime powers}

There are $\OO(X^{1/2})$ proper prime powers $p^k\le X$ with $k\ge2$.
For each such integer there are at most $R+1$ integers at most $X$ in the same residue class modulo $q$, and every von Mangoldt pair weight is at most $(\log X)^2$.
Thus the contribution of pairs in which at least one member is a proper prime power is
\begin{equation}\label{eq:pp-error}
 \OO\!\left(X^{1/2}(R+1)(\log X)^2\right)
 =o(X\log q),
\end{equation}
because $R\le\log q\asymp\log X$.
The remaining part of \cref{eq:weighted-pairs} is a sum over pairs of distinct primes in the same reduced residue class.  Since each weight is at most $(\log X)^2$,
\[
 C_2\gg\frac{X\log q}{(\log X)^2}.
\]
By \cref{eq:Xasymp,eq:logX}, there is an absolute $\alpha>0$ such that, for all sufficiently large $q$,
\begin{equation}\label{eq:C2}
 \boxed{C_2\ge\alpha\Phi.}
\end{equation}

\begin{checkpoint}[Second-moment closure]
The use of the variance theorem is pointwise in $q$, the centering transfer is $o(X\log q)$, the diagonal term cancels to $o(X\log q)$ at the selected cutoff, and the proper-prime-power contribution is also $o(X\log q)$.  No prime-pair conjecture remains.
\end{checkpoint}

\section{The third factorial moment from a uniform upper sieve}\label{sec:triple}

Every unordered triple $p_0<p_1<p_2\le X$ lying in one residue class modulo $q$ has a unique representation
\[
 p_1=p_0+rq,\qquad p_2=p_0+sq,
 \qquad 1\le r<s<R.
\]
Therefore
\begin{equation}\label{eq:C3-decomp}
 C_3\le\sum_{1\le r<s<R}T_q(r,s),
\end{equation}
where
\[
 T_q(r,s):=
 \#\{p\le X:p,\ p+rq,\ p+sq\text{ are prime}\}.
\]
The enlargement of the original range $p\le X-sq$ to $p\le X$ only increases the count.

For a prime $\ell$, let $\nu_\ell(r,s;q)$ be the number of distinct residues among
\[
 0,\ -rq,\ -sq\pmod\ell
\]
and define
\begin{equation}\label{eq:series}
 \Ss_q(r,s):=
 \prod_\ell
 \left(1-\frac{\nu_\ell(r,s;q)}{\ell}\right)
 \left(1-\frac1\ell\right)^{-3}.
\end{equation}

\subsection{Admissible and nonadmissible triples}

For the forms
\[
 p,\qquad p+rq,\qquad p+sq,
\]
Dubbe's nondegeneracy determinant is
\[
 rs(s-r)q^3\ne0.
\]
If the triple is admissible, \cref{thm:Dubbe} gives, uniformly in $q,r,s$,
\begin{equation}\label{eq:triple-sieve}
 T_q(r,s)\ll \Ss_q(r,s)\frac{X}{(\log X)^3}.
\end{equation}

If the triple is not admissible, some prime $\ell$ divides at least one of the three values for every integer $p$.  Since three residue classes cannot cover all classes modulo a prime $\ell>3$, necessarily $\ell\le3$.  If all three values are prime, one of them must equal $\ell$.  Each of the equations
\[
 p=\ell,\qquad p+rq=\ell,\qquad p+sq=\ell
\]
has at most one integer solution.  Hence
\begin{equation}\label{eq:nonadm}
 T_q(r,s)\le3
\end{equation}
for every nonadmissible pair $(r,s)$.

\subsection{A uniform singular-series majorant}

If $\ell\mid q$, then $\nu_\ell=1$, so the local factor is exactly
\[
 \left(1-\frac1\ell\right)^{-2}.
\]
Thus the product of the local factors at primes dividing $q$ is
\begin{equation}\label{eq:qfactor}
 \prod_{\ell\mid q}\left(1-\frac1\ell\right)^{-2}
 =\left(\frac q\Phi\right)^2.
\end{equation}

Now suppose $\ell\nmid q$ and $\ell\ge5$.
If $\ell\nmid rs(s-r)$, the three roots are distinct and $\nu_\ell=3$.  The local factor is less than $1$, because
\[
 (\ell-1)^3-\ell^2(\ell-3)=3\ell-1>0.
\]
If $\ell\mid rs(s-r)$, then $\nu_\ell\in\{1,2\}$, and the local factor is at most
\[
 \left(1-\frac1\ell\right)^{-2}\le1+\frac4\ell.
\]
The last inequality follows from
\[
 (\ell+4)(\ell-1)^2-\ell^3=2\ell^2-7\ell+4>0
 \qquad(\ell\ge5).
\]
The finitely many local possibilities at $2$ and $3$ contribute an absolute factor at most $16$.
Define
\[
 F(n):=\prod_{\substack{\ell\mid n\\\ell\ge5}}
 \left(1+\frac4\ell\right).
\]
Then
\begin{equation}\label{eq:series-majorant}
 \boxed{
 \Ss_q(r,s)\le16\left(\frac q\Phi\right)^2
 F(r)F(s)F(s-r).}
\end{equation}

\subsection{Average of the majorant}

For every prime $\ell\ge5$,
\[
 \left(1+\frac4\ell\right)^3\le1+\frac{52}{\ell}.
\]
Consequently
\begin{align}
 F(n)^3
 &\le\prod_{\ell\mid n}\left(1+\frac{52}{\ell}\right)\\
 &=\sum_{d\mid n}\frac{\mu^2(d)52^{\om(d)}}d.
\end{align}
Therefore
\begin{align}
 \sum_{n<R}F(n)^3
 &\le R\sum_{d\ge1}\frac{\mu^2(d)52^{\om(d)}}{d^2}\\
 &=R\prod_\ell\left(1+\frac{52}{\ell^2}\right)\\
 &\ll R.\label{eq:Fmean}
\end{align}
The Euler product converges absolutely.
Applying H\"older's inequality over $1\le r<s<R$ and using \cref{eq:Fmean} in the three variables $r,s,s-r$ gives
\begin{equation}\label{eq:Holder}
 \sum_{1\le r<s<R}F(r)F(s)F(s-r)\ll R^2.
\end{equation}
Combining \cref{eq:series-majorant,eq:Holder},
\begin{equation}\label{eq:series-average}
 \sum_{1\le r<s<R}\Ss_q(r,s)
 \ll R^2\left(\frac q\Phi\right)^2.
\end{equation}
This estimate is uniform for all moduli, including primorial and other very smooth moduli.

\subsection{Completion of the third-moment bound}

By \cref{eq:C3-decomp,eq:triple-sieve,eq:nonadm,eq:series-average},
\begin{equation}\label{eq:C3-pre}
 C_3\ll
 \frac{X}{(\log X)^3}R^2\left(\frac q\Phi\right)^2+R^2.
\end{equation}
Since
\[
 R\le\frac{\Phi\log q}{q},
\]
we have
\[
 R^2\left(\frac q\Phi\right)^2\le(\log q)^2.
\]
Hence the first term in \cref{eq:C3-pre} is
\[
 \ll\frac{X(\log q)^2}{(\log X)^3}\ll\Phi
\]
by \cref{eq:Xasymp,eq:logX}.  Also $R^2\le(\log q)^2=o(\Phi)$, using $\Phi\gg q/\log\log q$.  Thus there is an absolute $\beta>0$ such that, for all sufficiently large $q$,
\begin{equation}\label{eq:C3}
 \boxed{C_3\le\beta\Phi.}
\end{equation}

\begin{checkpoint}[Third-moment closure]
The upper sieve is uniform in the varying shifts, nonadmissible patterns contribute only $\OO(R^2)$, and the entire smooth-modulus factor is exactly $(q/\ph(q))^2$.  The latter cancels against the number of shift pairs.  No hidden $\log\log q$ loss remains.
\end{checkpoint}

\section{Finite extraction of repeated and empty classes}

Choose a fixed integer $T\ge3$ so large that
\begin{equation}\label{eq:Tchoice}
 \frac{3\beta}{T-1}\le\frac\alpha2.
\end{equation}
For every integer $k>T$,
\[
 \binom{k}{2}=\frac{3}{k-2}\binom{k}{3}
 \le\frac{3}{T-1}\binom{k}{3}.
\]
It follows from \cref{eq:C3,eq:Tchoice} that the contribution of the classes with $N_a>T$ to $C_2$ is at most $\alpha\Phi/2$.  By \cref{eq:C2},
\[
 \sum_{2\le N_a\le T}\binom{N_a}{2}\ge\frac\alpha2\Phi.
\]
Each summand is at most $\binom T2$, so
\begin{equation}\label{eq:D2}
 D_2:=\#\{a\pmod q:(a,q)=1,\ N_a\ge2\}
 \ge\delta\Phi,
 \qquad
 \delta:=\frac{\alpha}{2\binom T2}>0.
\end{equation}

Let
\[
 M:=\sum_{a\bmod q}^{*}N_a.
\]
For all sufficiently large $q$, $X\ge q$, so every prime divisor of $q$ is at most $X$.  Therefore
\[
 M=\pi(X)-\om(q).
\]
By the prime number theorem and \cref{eq:XlogX},
\begin{equation}\label{eq:M}
 M=\Phi+o(\Phi).
\end{equation}
Let $O(X)$ be the number of occupied reduced residue classes and let
\[
 U(X):=\Phi-O(X)
\]
be the number of empty reduced classes.  Every occupied class uses one prime incidence, and every class counted by $D_2$ uses at least one additional incidence.  Hence
\[
 M\ge O(X)+D_2.
\]
Using \cref{eq:D2,eq:M},
\begin{equation}\label{eq:UX}
 U(X)\ge\Phi-M+D_2\ge(\delta-o(1))\Phi.
\end{equation}

This is the finite implication checked by the accompanying Lean development, conditional on the two moment bounds and the total-incidence estimate.

\section{Enlarging the cutoff}

Let
\[
 Y:=(1+c)\Phi\log q,
\]
where $c>0$ is fixed.  From \cref{eq:Xasymp} and the prime number theorem,
\begin{equation}\label{eq:increment}
 \pi(Y)-\pi(X)=c\Phi+o(\Phi).
\end{equation}
Indeed, $Y/X\to1+c$, $\log X\sim\log q$, and the floor discrepancy in \cref{eq:floor} contributes only $o(\Phi)$ primes.

Each newly admitted prime can fill at most one class that was empty at $X$.  Thus
\[
 U(Y)\ge U(X)-\bigl(\pi(Y)-\pi(X)\bigr)
 \ge(\delta-c-o(1))\Phi.
\]
Choose
\[
 c:=\frac\delta4.
\]
Then, for all sufficiently large $q$,
\begin{equation}\label{eq:UY}
 U(Y)\ge\frac\delta2\Phi.
\end{equation}
If a reduced residue class contains no prime at most $\lfloor Y\rfloor$, its least congruent prime, being an integer, is strictly greater than $Y$.  We have proved the following.

\begin{theorem}[Candidate resolution of Erd\H{o}s Problem 971]\label{thm:main}
There exist absolute constants $c,C>0$ and $q_0$ such that, for every integer $q\ge q_0$,
\[
 \#\{a\pmod q:(a,q)=1,\ p(a,q)>(1+c)\ph(q)\log q\}
 \ge C\ph(q).
\]
One may take $c=\delta/4$ and $C=\delta/2$, with $\delta$ defined in \cref{eq:D2}.
\end{theorem}

\section*{Lean status and exact claim}
\addcontentsline{toc}{section}{Lean status and exact claim}

The existing Lean 4 development proves the following finite implication without project-local axioms or admitted proof terms:

\begin{quote}
If the second factorial moment has a positive linear lower bound, the third factorial moment has a linear upper bound, the total mass at the base cutoff is at most $(1+o(1))$ times the number of classes, and the cutoff extension adds sufficiently little mass, then a positive linear number of reduced residue classes have least congruent prime beyond the enlarged cutoff.
\end{quote}

Its strongest arithmetic specialization is named
\[
 \texttt{prime\_moment\_method\_to\_least\_prime\_classes}.
\]
It also proves
\[
 N_a(X)=0\quad\Longleftrightarrow\quad X<p(a,q)
\]
for reduced classes, using mathlib's formalized Dirichlet theorem.

The present paper supplies the missing informal analytic instantiation:
\[
 C_2\ge\alpha\Phi,\qquad C_3\le\beta\Phi,
 \qquad M=\Phi+o(\Phi),
 \qquad \pi(Y)-\pi(X)=c\Phi+o(\Phi).
\]
However, importing these as assumptions into Lean is not the same as proving them in Lean.  A genuinely complete formal verification would require substantial new mathlib developments for:

\begin{enumerate}[leftmargin=2em]
 \item Chebyshev functions and their second weighted moment in the required asymptotic form;
 \item the pointwise Friedlander--Goldston variance theorem;
 \item the explicit or abstract three-dimensional upper-bound sieve;
 \item Euler-product singular-series estimates and the H\"older average above;
 \item the PNT and totient estimates assembled with all uniform quantifiers and floors.
\end{enumerate}



\section*{Conclusion}
\addcontentsline{toc}{section}{Conclusion}

The earlier bottlenecks were not instances of a remaining parity obstruction.  The necessary pair lower bound is already encoded in the known pointwise variance lower bound, and the concentration problem is controlled by a standard three-tuple upper sieve.  After the singular series is averaged with its exact $q$-local factors, the argument closes for every sufficiently large modulus.

This provides a candidate affirmative solution of Erd\H{o}s Problem 971.  The next step is independent review by an analytic number theorist and, if the proof survives that review, circulation as a conventional preprint or journal submission.  The lack of end-to-end Lean formalization should be stated explicitly and should not be conflated with a gap in the informal analytic proof.

\begin{thebibliography}{99}

\bibitem{Dubbe2024}
T. Dubbe,
\emph{Explicit bounds for prime $k$-tuplets},
arXiv:2409.04261, 2024.
\url{https://arxiv.org/abs/2409.04261}

\bibitem{Erdos1950}
P. Erd\H{o}s,
\emph{On some applications of Brun's method},
Acta Sci. Math. (Szeged) \textbf{13} (1950), 57--63.

\bibitem{Fiorilli2015}
D. Fiorilli,
\emph{The distribution of the variance of primes in arithmetic progressions},
Int. Math. Res. Not. IMRN 2015, no. 12, 4421--4448;
arXiv:1301.5663.
\url{https://arxiv.org/abs/1301.5663}

\bibitem{FG1996}
J. B. Friedlander and D. A. Goldston,
\emph{Variance of distribution of primes in residue classes},
Q. J. Math. \textbf{47} (1996), 313--336.
\url{https://doi.org/10.1093/qmath/47.3.313}

\bibitem{Leung2025}
S.-K. Leung,
\emph{Moments of primes in progressions to a large modulus},
Forum Math. (2025), published online;
arXiv:2402.07941.
\url{https://arxiv.org/abs/2402.07941}

\end{thebibliography}

\end{document}
